%; whizzy section -pdf xpdf -latex ./whizzypdfptex.sh
% latex beamer presentation.
% platex, latex-beamer でコンパイルすることを想定。 

%     Tokyo Debian Meeting resources
%     Copyright (C) 2008 Junichi Uekawa

%     This program is free software; you can redistribute it and/or modify
%     it under the terms of the GNU General Public License as published by
%     the Free Software Foundation; either version 2 of the License, or
%     (at your option) any later version.

%     This program is distributed in the hope that it will be useful,
%     but WITHOUT ANY WARRANTY; without even the implied warranty of
%     MERCHANTABILITY or FITNESS FOR A PARTICULAR PURPOSE.  See the
%     GNU General Public License for more details.

%     You should have received a copy of the GNU General Public License
%     along with this program; if not, write to the Free Software
%     Foundation, Inc., 51 Franklin St, Fifth Floor, Boston, MA  02110-1301 USA

\documentclass[cjk,dvipdfmx,12pt]{beamer}
\usetheme{Tokyo}
\usepackage{monthlypresentation}

%  preview (shell-command (concat "evince " (replace-regexp-in-string "tex$" "pdf"(buffer-file-name)) "&"))
%  presentation (shell-command (concat "xpdf -fullscreen " (replace-regexp-in-string "tex$" "pdf"(buffer-file-name)) "&"))
%  presentation (shell-command (concat "evince " (replace-regexp-in-string "tex$" "pdf"(buffer-file-name)) "&"))

%http://www.naney.org/diki/dk/hyperref.html
%日本語EUC系環境の時
\AtBeginDvi{\special{pdf:tounicode EUC-UCS2}}
%シフトJIS系環境の時
%\AtBeginDvi{\special{pdf:tounicode 90ms-RKSJ-UCS2}}

\title{東京エリア Debian 勉強会}
\subtitle{資料}
\author{上川 純一 dancer@debian.or.jp\\IRC nick: dancerj}
\date{2008年11月15日}
\logo{\includegraphics[width=8cm]{image200607/openlogo-light.eps}}

\begin{document}

\frame{\titlepage{}}


\section{Intro}

\emtext{設営準備にご協力ください}

\section{}
\begin{frame}
 \frametitle{Agenda}
\begin{minipage}[t]{0.45\hsize}
  \begin{itemize}
  \item 注意事項
	\begin{itemize}
	 \item 飲食禁止
	 \item 政治/宗教/営利活動禁止
	 \item ustream にて試験ストリーミング中
	\end{itemize}
  \item 最近のDebian関連のイベント
	\begin{itemize}
	 \item 前回の勉強会、Debconf
	\end{itemize}
 \end{itemize}
\end{minipage} 
\begin{minipage}[t]{0.45\hsize}
 \begin{itemize}
  \item 「その場で勉強会資料を作成しちゃえ」 Debian を使った LaTeX 原稿作成合宿
 \end{itemize}
\end{minipage}
\end{frame}

\section{最近}

\begin{frame}
 \frametitle{9月}
\begin{minipage}[t]{0.45\hsize}
  \begin{itemize}
  \item 注意事項
	\begin{itemize}
	 \item 飲食禁止
	 \item 政治/宗教/営利活動禁止
	\end{itemize}
  \item quiz
  \item 最近のDebian関連のイベント
	\begin{itemize}
	 \item 前回の勉強会 
	\end{itemize}
 \end{itemize}
\end{minipage} 
\begin{minipage}[t]{0.45\hsize}
 \begin{itemize}
  \item Po4a でドキュメント翻訳の保守を楽にしよう
  \item 【でびあん】Debian パッケージメンテナというお仕事【現在募集中】
 \end{itemize}
\end{minipage}
\end{frame}


\section{DWN quiz}
\begin{frame}{Debian 常識クイズ}

Debian の常識、もちろん知ってますよね?
知らないなんて恥ずかしくて、知らないとは言えないあんなことやこんなこと、
みんなで確認してみましょう。

今回の出題範囲は\url{debian-devel-announce@lists.debian.org} に投稿された
内容とDebian Project Newsからです。

\end{frame}

\subsection{問題}


 \santaku
 {BTS　500000番はどんなバグ報告だったでしょうか}
 {SPAMだった}
 {cdbs に defoma 対応させるためのパッチ}
 {Lenny リリースが遅れているというバグ報告}
 {B}

 \santaku
 {10月25日から30日まで行われた BTS の景品は何でしょう？}
 {DPLからのキス}
 {次回のDebconf無料チケット}
 {おいしいクッキー}
 {C}
 
 \santaku
 {screenshots.debian.net は何をまとめたサイトでしょうか？}
 {デスクトップのスクリーンショット}
 {ソフトウェアのスクリーンショット}
 {Debian Developerのプライベート写真集}
 {B}

 \santaku
 {Joerg Jaspertがプロポーザルを出したDebian membershipとは？}
 {「Debian開発者」の現在の定義を変えるための提案}
 {Debian開発者の家族に関する提案}
 {Debianからforkしたディストリに関する提案}
 {A}
　% http://lists.debian.org/debian-devel-announce/2008/10/msg00005.html

 \santaku
 {Debian Installer team から出たアナウンスは何でしょう？}
 {ごっめーん！インストーラー遅れちゃった。}
 {Debian Installer lenny RC1 出たよ}
 {RC1 は飛ばして リリースする予定です。}
 {A}


\section{最近のDebian関連のイベント}
\emtext{Debconf8}
\emtext{OSC}


\section{事前課題紹介}
\emtext{事前課題の紹介}
% pre work home work

\begin{frame}{事前課題}

今回の事前課題は今回の勉強会の準備として、前日までに Debian GNU/Linux
(sid) のインストールされたノートパソコンを用意していただき、\LaTeX 環境を
構築し、Debian 勉強会資料が作成できるまでを確認してもらいました。これらの
準備で発生した問題などを {\bf \LaTeX 合宿準備奮闘記}として提出していただ
きました。

\end{frame}


% (query-replace "\\subsection" "\\end{frame}\\begin{frame}")
% (query-replace "\\subsubsection" "\\textbf")


% (query-replace "\\subsection" "\\end{frame}\\begin{frame}")
% (query-replace "\\subsubsection" "\\textbf")


\emtext{2008年計画}

\begin{frame}{2008年計画}

{\scriptsize
\begin{enumerate}
 \item 新年会「気合を入れる」
 \item Open Source Conference Tokyo (3/1)
 \item データだけのパッケージを作成してみる、
       ライセンスの考え方 (David Smith)
 \item バイナリ一つのパッケージを作成してみる (吉田@板橋)\\
       バージョン管理ツールを使いDebianパッケージを管理する(git)\\
       アップストリームの扱い(svn/git/cvs)(岩松 信洋さん)
 \item バイナリの分けたパッケージの作成。(前田さん)\\
       バイナリの分け方の考え方、アップグレードなどの運用とか。
 \item パッケージ作成(dpatch/debhelperで作成するパッケージ)(小林儀匡さん)\\
       man の書き方(roff or docbook)(でんさん)\\
       Open Source Conference Hokkaido
 \item パッケージ作成(kernel patch、kernel module)
       、Debconf発表練習
 \item Debconf アルゼンチン、Debian温泉

 \item Open Source Conference Tokyo/Fall、
       デーモン系のパッケージの作成、latex、 emacs-lisp、フォントパッケージ
 \item パッケージの cross-compile の方法、amd64 上で i386 のパッケージと
       か、OSC-Fall報告会、Debconf報告会
 \item 国際化 po-debconf / po化 / DDTP
 \item 忘年会
\end{enumerate}
}
\end{frame}

\emtext{LaTeX作成合宿}

\emtext{次回の勉強会}
\begin{frame}{次回の勉強会}
次回は忘年会です。
\end{frame}

\begin{frame}{宴会場所}
\begin{itemize}
 \item 宴会場所\\
       本日の宴会は「はなの舞」です。\\
       参加者はフロアに集合し、全員で移動しましょう。
 \item 片付け\\
       部屋を片付けるのにご協力ください。
\end{itemize}
\end{frame}

\end{document}

;;; Local Variables: ***
;;; outline-regexp: "\\([ 	]*\\\\\\(documentstyle\\|documentclass\\|emtext\\|section\\|begin{frame}\\)\\*?[ 	]*[[{]\\|[]+\\)" ***
;;; End: ***
